\begin{abstract}
We give a fault-tolerant quantum computation construction that combines a geometrically simple topological code with a local, fixed-rule error-correction procedure.  In particular, our construction uses the four-dimensional toric code and applies the same constant-depth sweep-rule correction cycle throughout the encoded computation.  It tolerates sufficiently weak local-stochastic Pauli noise in the circuit model, assuming noiseless classical processing with bounded depth per time step.  Unlike approaches based on hierarchical correction schedules, the recovery rule itself does not change with time: when the simulated computation is idle, the encoded circuit repeatedly executes the same correction cycle.  We do not claim a constant asymptotic overhead in the total circuit depth.

Our analysis adapts the explanation-tree method underlying proofs of Toom's theorem to the evolution of syndromes in topological quantum codes.  It also establishes two decoding results of independent interest.  First, we give an explicit implementation of the sweep rule for cubical $(d,d')$-toric codes with $d',d-d'\geq 2$, resolving the local-correction existence question posed by Kubica and Preskill.  Second, we prove that the \emph{ideal} (noiseless) sweep decoder corrects sufficiently weak local-stochastic errors, extending earlier guarantees for independently distributed errors: on a torus of side length $n$, $n^{d+1}$ ideal sweep cycles recover the encoded state except with probability $O(\operatorname{poly}(n))p^{\Omega(n)}$.  Thus the four-dimensional toric code provides a particularly simple setting in which a time-independent local correction rule admits both a fault-tolerance analysis under repeated noisy execution and an ideal-decoding theorem under correlated local-stochastic noise.
\end{abstract}
